\documentclass[10pt]{article}

\usepackage[utf8]{inputenc}

\usepackage[T1]{fontenc}

\usepackage{amsmath}

\usepackage{amsfonts}

\usepackage{amssymb}

\usepackage[version=4]{mhchem}

\usepackage{stmaryrd}

\usepackage{graphicx}

\usepackage[export]{adjustbox}

\graphicspath{ {./images/} }



\begin{document}

минимальный, т. к. он разлагается на смыслоразличительные сегменты «м» и «а», к-рые уже не могут быть разложены, т. е. они являются минимальными. Из всевозможных признаков, к-рыми могут характеризоваться звуки языка, в каждом конкретном языке лишь нек-рые являются различительными, и совокупность различительных признаков меняется от языка к языку. Так, долгота гласного, не имеющая различительного значения в русском языке, имеет его в латышском (напр., pile - капля, a pile - утка) и в ряде других языков. Предпринимались попытки формальной трактовки понятия фонемы с привлечением простейших математич. средств (см., напр., [4], [5]), однако достаточно полной формальной теории фонемы пока (1984) нет. С точки зрения различительных признаков в С. Л. рассматриваются и минимальные значимые единицы языка - м о р ф е мы. Напр., в русском слове «столику» выделяются: корневая морфема «стол»-, словообразовательная морфема - «ик» - и словоизменительная морфема - «у», несущая различительные признаки единственного числа и дательного падежа.



В С. л. разработаны т. н. дескриптивные процедуры исследования языка (см. [6], [7]), основанные на работе с информантом - носителем языка, к-рому исследователь задает вопросы типа «Правильно ли данное выражение?» и “Одинаковы или различны смыслы двух данных выражений?». Многие используемые в математич. лингвистике конструкции представляют собой, по существу, формализацию таких процедур (см. Аналитическая модель языка). Разработаны строго формальные, по существу математические, способы описания структуры предложения (см. Синтаксическая структура), позволившие изучить ряд важных проблем теории синтаксиса. Развивается с т р у к т у на я с е м а нт и к а, изучающая структуру отношений между смыслом языковых выражений и их формой и между смыслами различных выражений (см., напр., [10], [11]).



Развитие идей С. л. привело к выработке представления о языке как о механизме, служащем для порождения речевых выражений или для преобразования означаемых (смыслов) в означающие (тексты) и обратно (см. [8]-[10]). Это представление легло в основу теории формальных грамматик.



Лит.: [1] Д е С о с с ю $\mathrm{p}$ Ф., Труды по языкознанию, пер. с франц., М., 1977; [2] Т р у б е ц к о й Н. С., Основы фоно-



\includegraphics[max width=\textwidth, center]{2023_07_09_69b89c7e923dc9280051g-1(1)}

230 ; [4] У с пा е н с к и й В. А., «Вопросы языкознания», 1964 , № 6, с. 39-53; [5] Р е в 3 й н И. И., Структура языка как моделирующей системы, М., 1978; [6] Б л у м ф и л д Л., Язык, пер. с англ., M., 1968; [7] H a r r i s Z. S., Structural linguistics, Chi., 1961; ' [8] Х о м с к и й Н., в кн.: Новое в Јингвистике, в. 2, М., 1962, с. 412-527; [9] е г о же, Аспек-



\includegraphics[max width=\textwidth, center]{2023_07_09_69b89c7e923dc9280051g-1}

ч у к И. А., Опыт теории лингвистических моделей «Смыслњтекст», М., 1974; [11] А п р е с я н Ю. Д., Лексическая семантика, М., 1974; [12] е г о ж е, Идеи и методы современной структурной лингвистики, М., 1966. А. В. Гладкий.



СТРУКТУРНО УПОРЯДОЧЕННАЯ ГРУППА, p eше тотно упо р ядоченная г р уп п а, $l$-г р у п п а,- группа $G$, на множестве элементов к-рой задано отношение частичного порядка $\leqslant$, обладающее свойствами: 1) $G$ - решетка относительно $\leqslant$, т. е. для любых $x, y \in G$ существуют элементы $x \wedge y, x \vee y$ такие, что $x \wedge y \leqslant x, y$ и $x \vee y \geqslant x, y$; для любого $z \in G, z \leqslant x, y$ выполнено $z \leqslant x \wedge y$, и для любого $t \in G$ и $x, y \leqslant t$ выполнено $x \vee y \leqslant t ; 2)$ для любых $a, b, x, y \in G$ неравенство $a \leqslant b$ влечет за собой $x a y \leqslant x b y$. Эквивалентным образом С. у. г. может быть определена как алгебраич. система в сигнатуре $\langle\cdot,-1, e, \vee, \wedge\rangle$, удовлетворяющая аксиомам: 3$)\langle G, \cdot,-1, e\rangle$ - грушпа; 4) $\langle G, \vee, \wedge\rangle-$ решетка; 5) $x(y \vee z) t=x y t \vee x z t \quad$ и $x(y \wedge z) t=x y t \wedge x z t$ для любых $x, y, z, \quad t \in G$.



Решетка элементов С. у. г. дистрибутивна. М о д ул е м (соответственно по ло и и е ль н $о$ й и о т- р иц а т е ль ной ч а с тью) ә ле мен т а $x$ наз. элемент $|x|=x \vee x^{-1}$ (соответственно $x^{+=} x \vee e$ и $x^{-}=$ $=x \wedge e$ ). В С. У. г. верны соотношения:



\[

\begin{gathered}

x=x^{+} x^{-},|x|^{-1} \leqslant x \leqslant|x|, \\

|x|=x^{+}\left(x^{-}\right)^{-1}, x^{+} \wedge\left(x^{-}\right)^{-1}=e, \\

(x \vee y)^{-1}=x^{-1} \wedge y^{-1},(x \wedge y)^{-1}=x^{-1} \vee y^{-1} .

\end{gathered}

\]



Әлементы $x$ и $y$ наз. о р т о г о на ль ны м и, если $|x| \vee|y|=e$. Ортогональные әлементы перестановочны.



Подмножество $H \quad l$-группы $G$ наз. $l$-п о д г р у IIп о й, если $H$ - подгруппа и подрешетка в $G$; $l$-подгруппа $H$ наз. $l$-и д е а л о м С. у. г. $G$, если она нормальна и выпукла в $G$. Множество $l$-подгрупп С. У. г. образует подрешетку решетки всех ее подгрупп. $\mathrm{Pe}$ шетка $l$-идеалов С. У. г. дистрибутивна. $l$ - г о м о м о рф и з м о м $l$-группы $G$ в $l$-группу $H$ наз. гомоморфизм $\varphi$ группы $G$ в группу $H$ такой, что



\[

\varphi(x \vee y)=\varphi(x) \vee \varphi(y), \varphi(x \wedge y)=\varphi(x) \wedge \varphi(y) .

\]



Ядрами $l$-гомоморфизмов являются в точности $l$-идеалы $l$-групп. Если $G$ есть $l$-группа, $M \subset G$, то множество $M^{\perp}=\{x \in G|\quad| x|\wedge| m \mid$-е для всякого $m \in M\}$ является выпуклой $l$-подгруппой в $G$.



Группа $A(L)$ взаимно однозначных сохраняющих порядок отображений линейно упорядоченного множест ва $L$ на себя есть $l$-группа (если для $f, g \in A(L)$ положить $f \leqslant g$ тогда и только тогда, когда $f(\alpha) \leqslant g(\alpha)$ для любого $\alpha \in L)$. Всякая $l$-группа $l$-изоморфна $l$-подгруппе С. у. г. $A(L)$ для нек-рого подходящего множества $L$.



Класс всех С. у. г. является многообразием сигнатуры $\left\langle\cdot,^{-1}, e, \wedge, \vee\right\rangle$. Важнейшее его подмногообразиекласс С. У. г., аппроксимирующихся линейно упорядоченными группами (класс п р е д с т а в и м ы Г р у пा П).



Јит.: [1] Б и р к го ф Г., Теория структур, пер. с англ., М., 1952; [2] Ф у к с Л., Частично упорядоченные алгебраич.

системы, пер. с англ., М., 1965.



СТРӰКТУРНОЕ ПРОСТРАНСТВО к о л ь ц а - множество $\mathfrak{P}$ всех его примитивных идеалов с топологиеї, замкнутые множества в к-рой суть такие подмножества $C \subseteq \mathfrak{B}$, что $C$ содержит всякий идеал, содержащий пересечение всех идеалов из $C$ (ср. Зариского топология). С. П. кольца $R$ гомеоморфно С. п. факторкольца $R / J$, где $J$ - радикал Джекобсона. С. п. является $T_{0}$ пространством. Если все пріпитивные идеалы колыца максимальны, то С. п. оказывается $T_{1}$-пространством. С. п. кольца с единицей компактно. С. п. бирегулярного кольца (см. Регулярное кольцо) локально компактно и вполне несвязно. Оно используется для прсдставления бирегулярного кольца в виде кольца непрерывных функций с компактными носителями.



М., 1961. [1] Д ж е к о б с о н Н., Строение колец, пер. с англ., СТРУ КТУРНЫЙ ИЗОМОРФИЗМ - изоморфиям между решетками (структурами) двух однотипных алгебраич. систем.



Для групп исследовался вопрос о том, когда из С. и. двух групп следует их изоморфизм (см. [1]).



Лит.: [1] К у р о ш А. Г., Теория групп, 3 изд., М., 1967.



СТРУХАЛЯ ЧИСЛО - критерий подобия нестационарных движений жидкостей или газов. С. ч. характеризует одинаковость протекания процессов во времени:



\[

\mathrm{Sh}=l / v t=\omega l / v,

\]



где $v$ - характерная скорость течения, $l$ - характерный линейный размер, $t-$ характерный для нестационарного движения промежуток времени, $\omega-\mathbf{x}$ рактерная частота (иногда через Sh обозначают обратную величину $v t / l$ ).



Аналогичный критерий $H_{0}=v t / l$ в механических, тепловых и әлектромагнитных процессах наз. к р и т ери и м гомох ромности.





\end{document}